% Lemma 3.2.7 — minimal p-power in a principal local ideal.
% Thesis §3.2 (pp. 18–19), corrected per ERRATA E6.2 (standing hypotheses
% restated; false for split p) and E6.7 (shift variable renamed s -> a).

\begin{lemma}[Minimal $p$-power in a principal local ideal]
  \label{lem-3-2-7}
  \lean{QuadraticOrder.comap_map_span_tau_sub_eq_canonicalIdeal}
  \uses{notation, lem-3-2-2, lem-3-2-4}
  Let $d < 0$ be a discriminant ($d \equiv 0, 1 \pmod 4$), $d = D f^2$ with
  $D$ fundamental and $f \geq 1$ the conductor, and let
  $\mathcal{O} = \mathbb{Z}[\tau]$ with $\tau = (d + \sqrt{d})/2$ be the order
  of discriminant $d$ in $K = \mathbb{Q}(\sqrt{d})$, so that the minimal
  polynomial of $\tau$ is $g(x) = x^2 - dx + \frac{d^2 - d}{4}$, the
  conjugate is $\bar{\tau} = d - \tau$, and $N(\tau - a) = g(a) > 0$ for all
  $a \in \mathbb{Z}$. Let $p$ be a rational prime with $p \mid f$; then
  $g \bmod p$ is the square of a linear polynomial, so $\mathcal{O}$ has a
  unique prime ideal $\mathfrak{p}$ containing $p$, and the localization
  $\mathcal{O}_{\mathfrak{p}} \subseteq K$ satisfies
  $\mathcal{O}_{\mathfrak{p}} = \mathbb{Z}_{(p)}[\tau]$ by
  Lemma~\ref{lem-3-2-2}.

  Let $r \geq 0$ be an integer, let $a \in \mathbb{Z}$, and set
  \[
    \nu := v_p\bigl(N(\tau - a)\bigr) = v_p\bigl(g(a)\bigr) \geq 0 .
  \]
  Then the contraction to $\mathcal{O}$ of the principal ideal
  $p^r(\tau - a)\,\mathcal{O}_{\mathfrak{p}}$ is the canonical ideal of
  Lemma~\ref{lem-3-2-4} with parameters $k = r + \nu$, $m = 2r + \nu$,
  $A = a$:
  \[
    p^r(\tau - a)\,\mathcal{O}_{\mathfrak{p}} \cap \mathcal{O}
      = \bigl(p^{r+\nu},\; p^r(\tau - a)\bigr) .
  \]
  In particular, every nonzero rational number lying in
  $p^r(\tau - a)\,\mathcal{O}_{\mathfrak{p}}$ has $p$-adic valuation at least
  $r + \nu$, and the smallest positive rational integer (equivalently, the
  smallest power of $p$) contained in
  $p^r(\tau - a)\,\mathcal{O}_{\mathfrak{p}}$ is exactly $p^{r+\nu}$.
\end{lemma}

\begin{proof}
  \uses{notation, lem-3-2-2, lem-3-2-4}
  Since $d < 0$, $g(a) = (a - d/2)^2 - d/4$ is a positive rational integer,
  so we may write $g(a) = p^{\nu} u$ with $u \in \mathbb{Z}_{>0}$ and
  $p \nmid u$; note $u \in \mathbb{Z}_{(p)}^{\times}$.

  \emph{Step 1: quadratic relation.} The elements $\tau - a$ and
  $\bar{\tau} - a$ have sum $d - 2a$ and product $N(\tau - a) = g(a)$, so
  \begin{equation}
    (\tau - a)^2 = (d - 2a)(\tau - a) - g(a). \tag{1}
  \end{equation}

  \emph{Step 2: the local ring as a free module.} By
  Lemma~\ref{lem-3-2-2} and $p \mid f$ (which makes $p$ non-split),
  $\mathcal{O}_{\mathfrak{p}} = \mathbb{Z}_{(p)}[\tau]$. Since
  $\tau^2 = d\tau - \frac{d^2 - d}{4}$, we have
  $\mathbb{Z}_{(p)}[\tau] = \mathbb{Z}_{(p)} \oplus \mathbb{Z}_{(p)}\tau$,
  the sum being direct because $1, \tau$ are $\mathbb{Q}$-linearly
  independent ($\tau \notin \mathbb{Q}$ as $d < 0$). Likewise
  $\mathcal{O} = \mathbb{Z} \oplus \mathbb{Z}\tau$, and coordinates with
  respect to the $\mathbb{Q}$-basis $(1, \tau)$ of $K$ are unique.

  \emph{Step 3: normal form of the local ideal.} We claim
  \begin{equation}
    p^r(\tau - a)\,\mathcal{O}_{\mathfrak{p}}
      = \mathbb{Z}_{(p)}\, p^{r+\nu} + \mathbb{Z}_{(p)}\, p^r(\tau - a).
      \tag{2}
  \end{equation}
  For $\beta = e_0 + e_1\tau \in \mathcal{O}_{\mathfrak{p}}$
  ($e_0, e_1 \in \mathbb{Z}_{(p)}$), write
  $\beta = (e_0 + a e_1) + e_1(\tau - a)$; then by (1),
  \begin{equation}
    p^r(\tau - a)\beta
      = \lambda \cdot p^r(\tau - a) + \mu \cdot p^{r+\nu},
    \qquad
    \lambda := e_0 + (d - a)e_1, \quad \mu := -u e_1 . \tag{3}
  \end{equation}
  The $\mathbb{Z}_{(p)}$-linear map $(e_0, e_1) \mapsto (\lambda, \mu)$ has
  determinant $-u \in \mathbb{Z}_{(p)}^{\times}$, hence is a bijection of
  $\mathbb{Z}_{(p)}^2$. As the principal ideal consists exactly of the
  elements $p^r(\tau - a)\beta$ with $\beta \in \mathcal{O}_{\mathfrak{p}}$,
  claim (2) follows.

  \emph{Step 4: containment $\supseteq$.} Clearly
  $p^r(\tau - a)$ lies in both $\mathcal{O}$ and the local ideal. For the
  integer generator, $\bar{\tau} - a = (d - a) - \tau \in \mathcal{O}$ and
  $p^r(\tau - a)(\bar{\tau} - a) = p^r g(a) = p^{r+\nu} u$, so
  \[
    p^{r+\nu}
      = p^r(\tau - a) \cdot u^{-1}(\bar{\tau} - a)
      \in p^r(\tau - a)\,\mathcal{O}_{\mathfrak{p}},
  \]
  using $u^{-1} \in \mathbb{Z}_{(p)} \subseteq \mathcal{O}_{\mathfrak{p}}$.
  The contraction is an ideal of $\mathcal{O}$ containing both generators,
  hence contains $\bigl(p^{r+\nu}, p^r(\tau - a)\bigr)$.

  \emph{Step 5: containment $\subseteq$.} Let
  $\xi \in p^r(\tau - a)\,\mathcal{O}_{\mathfrak{p}} \cap \mathcal{O}$. By
  (2), $\xi = \mu\, p^{r+\nu} + \lambda\, p^r(\tau - a)$ with
  $\lambda, \mu \in \mathbb{Z}_{(p)}$, and in the basis $(1, \tau)$,
  \[
    \xi = \bigl(\mu\, p^{r+\nu} - a\lambda\, p^r\bigr) \cdot 1
        + \bigl(\lambda\, p^r\bigr) \cdot \tau .
  \]
  Since $\xi \in \mathcal{O} = \mathbb{Z} \oplus \mathbb{Z}\tau$ and
  coordinates are unique, $\lambda p^r \in \mathbb{Z}$ and
  $\mu p^{r+\nu} - a\lambda p^r \in \mathbb{Z}$. From
  $\lambda \in \mathbb{Z}_{(p)}$ we get $v_p(\lambda) \geq 0$, and from
  $\lambda p^r \in \mathbb{Z}$ we get $v_q(\lambda) \geq 0$ for every prime
  $q \neq p$; hence $\lambda \in \mathbb{Z}$. Then
  $a \lambda p^r \in \mathbb{Z}$ forces $\mu p^{r+\nu} \in \mathbb{Z}$, and
  the same argument gives $\mu \in \mathbb{Z}$. Thus $\xi$ is a
  $\mathbb{Z}$-linear (in particular $\mathcal{O}$-linear) combination of
  $p^{r+\nu}$ and $p^r(\tau - a)$, so
  $\xi \in \bigl(p^{r+\nu}, p^r(\tau - a)\bigr)$.

  Steps 4 and 5 prove the contraction identity. The parameters are
  admissible for Lemma~\ref{lem-3-2-4}: with $k = r + \nu$ and
  $m = 2r + \nu$ we have $m/2 \leq k \leq m$ and
  $g(a) \equiv 0 \pmod{p^{2k - m}}$ since $2k - m = \nu = v_p(g(a))$.

  \emph{Step 6: minimality.} Let
  $\rho \in p^r(\tau - a)\,\mathcal{O}_{\mathfrak{p}} \cap \mathbb{Q}$ be
  nonzero. By (2), $\rho = \mu\, p^{r+\nu} + \lambda\, p^r(\tau - a)$ with
  $\lambda, \mu \in \mathbb{Z}_{(p)}$; its $\tau$-coordinate $\lambda p^r$
  must vanish, so $\lambda = 0$, $\rho = \mu p^{r+\nu}$, and
  $v_p(\rho) = v_p(\mu) + r + \nu \geq r + \nu$. Conversely, every
  $n \in \mathbb{Z}$ with $v_p(n) \geq r + \nu$ lies in the local ideal,
  since $n = (n / p^{r+\nu}) \cdot p^{r+\nu}$ with
  $n / p^{r+\nu} \in \mathbb{Z}_{(p)}$ and $p^{r+\nu}$ is in the ideal by
  Step 4. Hence the rational integers in the local ideal are exactly those
  with $v_p \geq r + \nu$, and the smallest positive one is $p^{r+\nu}$.
\end{proof}

% Divergences from the thesis (flagged per ERRATA):
% (1) E6.7 — the thesis's shift variable s is renamed a, avoiding the
%     collision with Lemma 3.2.6's modulus exponent s; both lemmas are used
%     simultaneously in §3.3.
% (2) E6.2 — the standing hypothesis p | f (whence p non-split, unique
%     prime above p, and O_p = Z_(p)[tau]) is restated in the lemma; the
%     statement is FALSE for split p (e.g. d = -4, p = 5, a = 0, r = 0 at
%     the prime (5, tau - 1): the local ideal is the unit ideal, minimal
%     integer 1, not 5^{r+nu} = 5).
% (3) The thesis asserts only the minimal-positive-integer claim; the
%     contraction identity stated here is what the thesis's proof of
%     Lemma 3.2.8 actually uses and attributes to 3.2.7. The proof above
%     (normal form (2), in the (1, tau)-basis rather than the thesis's
%     (1, sqrt(d)/2)-basis with its rank-1 determinant argument) yields
%     both claims at once.
